Una empresa química fabrica tres tipos de productos químicos, P1, P2 y P3, cuyos beneficios unitarios son, respectivamente, \textbf{1, 2 y 4} unidades monetarias (u.m.) 

Para su producción, se consumen dos materias primas. En particular, se consumen \textbf{1} tonelada de materia prima M1 y \textbf{2} de materia prima M2 por cada tonelada de P1, \textbf{2} de M1 y \textbf{1} de M2 por cada tonelada de P2 y \textbf{2} de M1 y \textbf{4} de M2 por cada tonelada de P3.

Cada mes, se dispone de \textbf{200} y \textbf{500} toneladas de materias primas P1 y P2, respectivamente. 

Por cada tonelada de P2 se generan \textbf{2} toneladas de un subproducto peligroso. Por cada tonelada de P1 se genera \textbf{1} tonelada del mismo subproducto. A su vez, por cada tonelada de de P3 se consume \textbf{1} toneladas de dicho subproducto. La empresa es autosuficiente con este subproducto y no dispone de ninguna cantidad adicional comprada, de forma que todo lo que necesita para P2 se genera mediante la producción de P1 y P3.

Para obtener el programa de producción de máximo beneficio se dispone del siguiente modelo de programación lineal, donde $x_1$, $x_2$ y $x_3$ representan las toneladas mensuales producidas de P1, P2 y P3, respectivamente.
 
\begin{equation}
\begin{split}
\mbox{max. } z = x_1 + 2x_2 + 4x_3\\
s.a.:\\
 x_1 + 2x_2 + 2x_3 \leq 200\\
2x_1 +  x_2 + 4x_3 \leq 500\\
 x_1 - 2x_2 +  x_3 = 0\\
x_1,\,\,x_2,\,\,x_3\geq 0\\
\end{split}
\end{equation}

Se conoce la tabla correspondiente a la solución óptima.

\begin{center}
\begin{tabular}{c|c|ccccc|}
 & $z$ & $x_1$ & $x_2$ & $x_3$ & $h_1$ & $h_2$\\ 
      & $-1000/3$ & $-4/3$ & 0 & 0 & $-5/3$ & 0 \\ 
\hline
$x_2$ & 100/3  & $-1/6$ & 1 & 0 & $1/6$ & 0\\ 
$h_2$ & 200  & $-1/2$ & 0 & 0 & $-3/2$ & 1\\ 
$x_3$ & 200/3  & $2/3$ & 0 & 1 & $1/3$ & 0\\ 
\hline
\end{tabular}
\end{center}





\begin{enumerate}
    \item Interpretar la información que proporciona la tabla correspondiente a la solución óptima.

    El plan de producción que ofrece el mayor beneficio:
    \begin{itemize}
        \item Proporciona un beneficio total de 333.33 u.m.
        \item No se produce P1
        \item Se producen 33.33 toneladas de P2.
        \item Se producen 66.67 toneladas de P3.
        \item Se emplean por completo las 200 toneladas disponibles de M1.
        \item No se emplean 200 de las 500 toneladas disponibles de M2.
        
    \end{itemize}

    \item Dentro de qué rango de valores puede variar el beneficio del producto P2 sin que cambie la gama de productos que resulta más beneficioso producir.

    Es necesario realizar el análisis de sensibilidad de $c_2$.

    Para que la gama siga siendo la misma, la base óptima debe seguir siéndolo y eso pasa porque $V^B \leq 0$

    \begin{equation}
    \begin{split}
    V^B = c-c^BB^{-1}A = c-c^Bp^B =\\
    = \begin{pmatrix}
    1 & c_2 & 4 & 0 & 0\\
    \end{pmatrix}
     - \begin{pmatrix}
    c_2 & 0 & 4\\
    \end{pmatrix}
    \begin{pmatrix}
    -1/6 & 1 & 0 & 1/6 & 0\\
    -1/2 & 0 & 0 & -3/2 & 1\\
    2/3 & 0 & 1 & 1/3 & 0\\
    \end{pmatrix}
     = \\
     =\begin{pmatrix}
    1 & c_2 & 4 & 0 & 0\\
    \end{pmatrix} - 
     \begin{pmatrix}
    \frac{-c_2+16}{6}  &  c_2  & 4 & \frac{c_2+8}{6}  & 0 \\
    \end{pmatrix}=   \\  
     =\begin{pmatrix}
     \frac{-10+c_2}{6}& 0 & 0 & -\frac{c_2+8}{6} & 0\\
    \end{pmatrix} \leq 0
    \end{split}
    \end{equation}
        
    
    Por lo tanto, se debe cumplir $\frac{-10+c_2}{6} \leq 0$ y $-\frac{c_2+8}{6} \leq 0$, es decir, $c_2 \leq 10$ y $c_2 \geq -8$.

    Mientras el margen del producto P1 sea igual o superior a $-8$ (no de pérdidas superiores a 8 u.m.) y no sea mayor 10, la gama de productos es la que se indica en el apartado anterior.



    \item Ha surgido la necesidad de producir 1 tonelada de producto P1, ¿qué impacto tendría sobre el beneficio y sobre el resto de la producción?

    Producir una tonelada de P3 implicaría que $x_1 = 1$. En la solución óptima $x_1$ es una variable no básica y, por lo tanto, su valor es nulo. El impacto de aumentar una unidad $x_1$ queda reflejado en:

    \begin{itemize}
        \item El criterio del simplex, $V^B_3 = -4/3$. Por lo que al producir una tonelada adicional, el beneficio disminuiría en 4/3 u.m.
        \item Las tasas de sustitución dan cuenta de la modificación de los valores de las variables básicas. En este caso:
        \begin{itemize}
            \item $p^B_{11} = -1/6$, $x_2$ aumentaría 1/6, es decir, se producirían 166.67 kg más de $P_2$.
            \item $p^B_{21} = -1/2$, $h_2$ aumentaría en 1/2, es decir, se haría uso de 500 kg menos de $M_2$.
            \item $p^B_{31} = 2/3$, $x_3$ disminuiría en 2/3, es decir, se producirían 666.67 kg menos de producto $P_3$.
        \end{itemize}
    \end{itemize}

     
    
    \item ¿Qué impacto tendría sobre el beneficio la posibilidad de comprar y disponer de una tonelada de subproducto (diferente de las que se generan con la producción de P1 y P3)? Si la empresa dispone de ese producto adicional, dada la naturaleza peligrosa del mismo, debería consumirlo obligatoriamente.

    El hecho de que se disponga de una unidad más de subproducto sería equivalente a que la tercer restricción tuviera la forma:
    
    $$1x_1 - 2x_2 + 1x_3 = 1$$

    Es decir, el términos $b_3$ aumentaría una unidad. El precio sombra de la restricción 3 es igual al incremento de la función objetivo cuando se aumenta una unidad la disponibilidad del recurso correspondiente a esa restricción (parta derecha de la ecuación). Es decir, $\pi^B_3 = \Delta z|_{\Delta h_3 = 1}$

    Para calcular el precio sombra es necesario disponer de la inversa de la base:

    \begin{equation}
    \begin{split}
    B^{-1} =\begin{pmatrix}
    1/6 & 0 & -1/3\\
    -3/2 & 1 & -1\\
    1/3 & 0 & 1/3\\
    \end{pmatrix}
    \end{split}
    \end{equation}
    
    El precio sombra se calcula como:

    \begin{equation}
    \begin{split}
    \pi^B = c^BB^{-1} = \begin{pmatrix}
    2 & 0 & 4\\
    \end{pmatrix}
    \begin{pmatrix}
    1/6 & 0 & -1/3\\
    -3/2 & 1 & -1\\
    1/3 & 0 & 1/3\\
    \end{pmatrix}
     = \begin{pmatrix}
    5/3 & 0 & 2/3\\
    \end{pmatrix}
    \end{split}
    \end{equation}
    
    El precio sombra de la restricción 3 es 2/3, lo cual quiere decir que si se dispusiera de 1 tonelada del subproducto peligroso, la función objetivo aumentaría en 2/3 u.m.
    
\end{enumerate}






